\section{Experimental Configuration}
\label{sec:experiment}

This section specifies the complete experimental configuration required to reproduce our evaluation framework. Following the methodology established in Section~\ref{sec:methodology}, we provide concrete parameter values for dataset composition, LLM configurations, experimental tracks, classification settings and statistical analysis procedures. Detailed hyperparameter grids, prompt templates, and LLM configurations are provided in Appendices~\ref{sec:appendix}.

\subsection{Dataset and Sampling}
We employ the CEAS-08 spam email dataset~\cite{champa2024curated}, comprising 39,154 real-world email samples (21,842 spam, 17,312 legitimate) available in our repository\href{https://github.com/tisage/llm-synthetic-spam-benchmark}{\textsuperscript{[1]}}. The corpus is partitioned into $R = 20$ independent groups through stratified random sampling, with each group containing $N = 1{,}000$ samples at fixed spam ratio $\rho = 0.1$ (100 spam, 900 ham). This 10\% imbalance reflects realistic operational conditions. An 80-20 train-test split establishes a fixed held-out test set held constant across all experiments, eliminating test set variance as a confounding factor.

\subsection{LLM Models and Generation}

\textit{Ethics Note:} All synthetic spam emails are generated solely for research purposes to advance defensive cybersecurity capabilities. Complete prompt templates are provided in supplementary materials to support reproducibility.

We employ two state-of-the-art large language models: \textit{GPT-4.1-mini} (OpenAI) and \textit{Claude-3.5-Haiku} (Anthropic). These models offer optimal cost-performance trade-offs for realistic spam generation scenarios, providing sufficient context windows while maintaining low per-token costs economically viable for large-scale campaigns. Both models employ identical generation parameters: temperature = 0.7, max tokens = 1,000, top-p = 1.0. Complete API configurations, cost analysis, and parameter justifications are provided in Appendix~\ref{sec:appendix}.

We implement three systematically varied prompt strategies to explore generation intensity: \textbf{Original} (faithful reproduction with lexical variation), \textbf{Strong} (amplified spam indicators), and \textbf{Weak} (subtle spam with reduced explicit signals). Each prompt follows a consistent JSON-formatted template ensuring reliable parsing. Complete prompt templates are provided in supplementary materials to support experimental reproducibility. Table~\ref{tab:spam_examples_gpt} illustrates how prompt variation affects generation quality across the three strategies.


\vspace{0.5em}
\noindent\fbox{\parbox{0.95\columnwidth}{\small
\textbf{Prompt Template Structure:}

\texttt{You are an expert email content analyzer. Given the following spam email, generate a semantically similar but lexically distinct spam variant.}

\textbf{[\,Strategy-Specific Instructions\,]}

\texttt{Original Email:} \textbf{\{\,original\_spam\_content\,\}}

\texttt{Output Format (JSON):}\\
\texttt{\{\,"subject": "...", "body": "..."\,\}}
}}
\vspace{0.5em}


\begin{table*}[t]
\centering
\caption{Representative Synthetic Spam Examples (GPT-4.1-mini). Original prompts produce natural spam with realistic language patterns, strong prompts amplify urgency cues and explicit offers, while weak prompts generate subtle spam with professional tone.}
\label{tab:spam_examples_gpt}
\small
\begin{tabular}{p{0.12\textwidth}p{0.83\textwidth}}
\toprule
\textbf{Strategy} & \textbf{Generated Email} \\
\midrule
\textbf{Original} &
\textit{Subject:} Congratulations! You've Won \$5.5 Million \\
& \textit{Body:} Dear Winner, Your email has been randomly selected in our international lottery draw. You have won \$5,500,000.00! To claim your prize, contact our claims department immediately with your personal details. Act now - this offer expires in 48 hours! http://prize... \\
\midrule
\textbf{Strong} &
\textit{Subject:} URGENT: \$5.5M Lottery Prize - CLAIM NOW! \\
& \textit{Body:} IMMEDIATE ACTION REQUIRED! You are the GUARANTEED winner of \$5.5 MILLION DOLLARS! This is a LIMITED TIME OFFER - only 24 hours left! Don't let this life-changing opportunity slip away... Contact us RIGHT NOW or lose your prize forever! http://urgent... \\
\midrule
\textbf{Weak} &
\textit{Subject:} Notification Regarding Prize Selection \\
& \textit{Body:} Dear Recipient, We are pleased to inform you that your email address was included in our recent prize selection process. You may be eligible to receive a substantial financial award. For verification and processing, please provide your contact information at your earliest convenience... http://verify... \\
\bottomrule
\end{tabular}
\end{table*}

\subsection{Experimental Tracks}

Following the methodology framework (Section~\ref{sec:methodology}), we implement three experimental tracks with specific parameter configurations. All tracks employ cross-group mixing, no synthetic exposure, all three prompt strategies, and 20 independent replications. Table~\ref{tab:track_params} summarizes the key configuration differences.

\begin{table*}[t]
\centering
\caption{Track-Specific Parameters}
\label{tab:track_params}
\small
\begin{tabular}{lccc}
\toprule
\textbf{Parameter} & \textbf{Track 1} & \textbf{Track 2a} & \textbf{Track 2b} \\
\midrule
Training spam & Real + Synthetic & Real only & Real + Synthetic \\
Testing spam & Real & Synthetic & Synthetic \\
Controlled variable & Synthetic ratio & Training size & Synthetic ratio \\
 & $\theta$ & (real spam count) & $\theta$ (cross-model) \\
Variable values & 0\%, 10\%, ..., 100\% & 100\%, 150\%, 200\% & 0\%, 50\%, 100\% \\
\bottomrule
\end{tabular}
\end{table*}

\textbf{Track 1 (RQ1):} Varies synthetic proportion $\theta$ from 0\% to 100\% in 10\% increments to quantify performance degradation as real spam is replaced by synthetic spam during training. Training sets contain $(1-\theta) \times 100$ real spam + $\theta \times 100$ synthetic spam + 900 ham. Testing employs 100 real spam + 900 ham (fixed across all ratios).

\textbf{Track 2a (RQ2a):} Varies training data size to evaluate whether additional real spam improves zero-shot detection of LLM-generated synthetic spam. Training employs 100\%, 150\%, or 200\% real spam (100, 150, or 200 samples) with 900 ham. Testing evaluates detection of LLM-generated synthetic spam across both models and all prompt strategies.

\textbf{Track 2b (RQ2b):} Evaluates cross-model augmentation effectiveness at three synthetic ratios ($\theta = 0\%, 50\%, 100\%$). Training employs synthetic spam from one LLM while testing evaluates detection of synthetic spam from the opposite LLM. For example, training on Claude synthetic spam tests detection of GPT synthetic spam.

\subsection{Classification Configurations}

We employ two classifiers representing distinct learning paradigms: Support Vector Machine (SVM) and Random Forest (RF). All emails are transformed to TF-IDF feature vectors, with the vectorizer fitted exclusively on training data to prevent information leakage. Hyperparameters are optimized through grid search with 5-fold stratified cross-validation using F1-score as the optimization metric. Complete hyperparameter grids and TF-IDF configurations are provided in Appendices~\ref{sec:appendix}

\subsection{Evaluation Metrics}

We compute descriptive statistics (mean, standard deviation) across 20 independent groups for each configuration. Paired \textit{t}-tests compare each synthetic ratio against baseline with Benjamini-Hochberg FDR correction ($\alpha = 0.05$) for multiple comparisons. Cohen's $d$ effect sizes quantify practical importance beyond statistical significance. Linear regression models performance degradation as a function of synthetic proportion. F1-score serves as the primary metric, with supporting metrics including Accuracy, Precision, Recall, and AUC-ROC.

The complete evaluation framework comprises 4,080 training-evaluation units across all tracks (Track 1: 2,640; Track 2a: 720; Track 2b: 720), covering all combinations of LLM models, prompt strategies, classifiers, and replications. All experiments employ official LLM APIs (GPT-4.1-mini, Claude-3.5-Haiku) with total 12,000 synthetic samples. Classifier training executes on standard computing infrastructure (Intel Xeon 32-core CPU, 64 GB RAM). Complete cost breakdown, infrastructure specifications, and experimental results are provided in Appendix~\ref{sec:appendix} and our repository\href{https://github.com/tisage/llm-synthetic-spam-benchmark}{\textsuperscript{[1]}}.
